下表按环节如实记录生成式 AI 工具的使用方式、AI 的具体产出以及对应的人工核验方法。
本文件内容与仓库中的 \nolinkurl{AI_USAGE.md} 一致。

\begin{longtable}{@{}p{0.15\linewidth}p{0.40\linewidth}p{0.37\linewidth}@{}}
\toprule 环节 & AI 的具体产出 & 人工核验方式 \\ \midrule \endhead \bottomrule \endfoot
问题理解与歧义分析 &
读题后列出题面中未明确之处（区间开闭、冲突是否逐次比较、优先级如何形式化、问题 3“不限制平移幅度”指谁），
为每一处给出备选解释与取舍理由，形成建模决策记录（MDR-0001 至 MDR-0010）。 &
逐条对照题目原文与附件示例核对（如以 A001 的 $[35,40)$、$[100,105)$ 确认半开区间与推进规则）；
对无法由题面判定的歧义（问题 3），要求同时给出两种解释的结果而非单选。 \\
模型与算法设计 &
提出时频单元格的统一刻画、单元格 AtMostOne 编码、七级词典序目标、现任解保持的逐级搜索、
问题 3 的候选位置生成与上界链；给出命题与证明草稿。 &
逐条复核证明（命题~\ref{prop:cell}、\ref{prop:exact}、\ref{prop:weights}、\ref{prop:cut}、\ref{prop:area}、引理~\ref{lem:tiling}）；
对关键命题在程序中写成性质测试，由测试而非自然语言论证来担保。 \\
程序编写与调试 &
编写数据接入、契约校验、两种检测算法、消解与装填模型、独立校验器、图表与论文流水线的全部源程序与单元/性质/集成测试。 &
静态门禁（ruff、mypy）与 $53$ 项测试在云端全部通过；关键实现（校验器）要求与求解器代码分离，
不得复用同一函数。 \\
实验组织与缺陷诊断 &
组织云端分阶段运行；在发现“撤销 C”一级的结果劣于已知可行解后，设计对照实验定位原因
（逐级求解丢弃了热启动现任解），提出并实现现任解上界割。 &
以同模型、同固定等式、同起点的对照运行确认退化来自实现而非实例难度；
修复后重跑全部受影响阶段，并以独立校验器与微型实例穷举对照复核新解。 \\
数值核验 &
生成独立校验器、穷举对照、交叉验证与灵敏度分析的实验脚本与结果汇总。 &
人工核对表~\ref{tab:validation} 的每一行结论与各阶段 \nolinkurl{validation_report.json}；
核对结果工作簿的行数与统计表是否一致；对“已证明最优”与“已知最好”的措辞逐句把关。 \\
图表绘制 &
按色盲友好色板绘制全部 $\valFigureCount$ 张矢量图，生成 $\valTableCount$ 张三线表（同时输出 CSV）。 &
逐张查看图形并核对图注与正文陈述是否一致；核对表中数值与 \nolinkurl{numbers.json} 是否同源。 \\
论文文字组织 &
按 CUMCM 格式撰写与改写各章节文字，建立“每个数字只能来自阶段产物”的排版约束。 &
逐节通读并修改；抽查正文数字与结果文件、阶段产物的一致性；
确保未证明最优之处不出现“最优”字样，确保全文无身份信息。 \\
\end{longtable}

\section*{三、结果的责任归属}

本文的全部建模决策（歧义解释、附加假设、目标函数形式化、备选方案取舍）均由作者作出并记录在案；
AI 工具用于加速文献与方案的比较、程序实现与文字组织。
论文中的每一个数值都由云端运行的阶段产物自动注入，未经人工键入；
凡未能证明最优之处，论文均按“值 / 证明下界 / 状态”如实报告，不作最优性声明。
